\let\negmedspace\undefined
\let\negthickspace\undefined
\documentclass[journal,12pt,twocolumn]{IEEEtran}
\usepackage{cite}
\usepackage{amsmath,amssymb,amsfonts,amsthm}
\usepackage{algorithmic}
\usepackage{graphicx}
\usepackage{textcomp}
\usepackage{xcolor}
\usepackage{txfonts}
\usepackage{listings}
\usepackage{enumitem}
\usepackage{mathtools}
\usepackage{gensymb}
\usepackage{comment}
\usepackage[breaklinks=true]{hyperref}
\usepackage{tkz-euclide}
\usepackage{listings}
\usepackage{gvv}  
\usepackage{tikz}
\usepackage{circuitikz}
\usepackage{caption}
\def\inputGnumericTable{}              
\usepackage[latin1]{inputenc}          
\usepackage{color}                    
\usepackage{array}                    
\usepackage{longtable}                
\usepackage{calc}                     \usepackage{multirow}                 \usepackage{multicol}
\usepackage{hhline}                    
\usepackage{ifthen}                    
\usepackage{lscape}
\usepackage{amsmath}
\newtheorem{theorem}{Theorem}[section]
\newtheorem{problem}{Problem}
\newtheorem{proposition}{Proposition}[section]
\newtheorem{lemma}{Lemma}[section]
\newtheorem{corollary}[theorem]{Corollary}
\newtheorem{example}{Example}[section]
\newtheorem{definition}[problem]{Definition}
\newcommand{\BEQA}{\begin{eqnarray}}
\newcommand{\EEQA}{\end{eqnarray}}
\newcommand{\define}{\stackrel{\triangle}{=}}
\theoremstyle{remark}
\newtheorem{rem}{Remark}
\title{18-Definite Integrals and Applications of Integrals
\large{EE1030 : Matrix Theory}

Indian Institute of Technology Hyderabad
% }
}
\author{Sarvajith Guddety}

\begin{document}

\maketitle







\noindent\textbf{9.} Let \( S \) be the area of the region enclosed by \( y = e^{-x^2} \), \( y = 0 \), \( x = 0 \) and \( x = 1 \); then $\hfill \textbf{(2012)}$

\begin{multicols}{2}
\begin{enumerate}
\item  [(a)] \( S \geq \frac{1}{e} \)
\item [(b)]\( S \geq 1 - \frac{1}{e} \)
\columnbreak
\item  [(c)] \( S \leq \frac{1}{4} \left( 1 + \frac{1}{\sqrt{e}} \right) \)
\item  (d) \( S \leq \frac{1}{\sqrt{2}} + \frac{1}{\sqrt{e}} \left( 1 - \frac{1}{\sqrt{2}} \right) \)
\end{enumerate}
\end{multicols}

   

   \noindent\textbf{10.} The option(s) with the values of $a$ and $L$ that satisfy the following equation is(are)$ \hfill \textbf{(JEE Adv. 2015)}$

   \[
   \frac{\int_0^{4\pi} e^t \left( \sin^6(at) + \cos^4(at) \right) dt}{\int_0^{\pi} e^t \left( \sin^6(at) + \cos^4(at) \right) dt} = L?
   \]

   \begin{multicols}{2}
   \begin{enumerate}
   \item  [(a)] $a = 2$, $L = \frac{e^{4\pi} - 1}{e^\pi - 1}$
   \item [(b)] $a = 2$, $L = \frac{e^{4\pi} + 1}{e^\pi + 1}$
   \columnbreak
   \item  [(c)] $a = 4$, $L = \frac{e^{4\pi} - 1}{e^\pi - 1}$
   \item [(d)] $a = 4$, $L = \frac{e^{4\pi} + 1}{e^\pi + 1}$
   \end{enumerate}
   \end{multicols}

   \noindent\textbf{11.} Let $f(x) = 7\tan^8x + 7\tan^6x - 3\tan^4x - 3\tan^2x$ for all $x \in \left( -\frac{\pi}{2}, \frac{\pi}{2} \right)$. \\
   Then the correct expression(s) is(are) $\hfill \hfill \textbf{(JEE Adv. 2015)}$
    
    \begin{multicols}{2}
    \begin{enumerate}
    \item  [(a)] $\int_0^{\pi/4} x f(x) dx = \frac{1}{12}$ \item [(b)]$\int_0^{\pi/4} f(x) dx = 0$
    \columnbreak
    \item  [(c)] $\int_0^{\pi/4} x f(x) dx = \frac{1}{6}$
    \item  [(d)] $\int_0^{\pi/4} f(x) dx = 1$ \end{enumerate}
    \end{multicols}

     

     \noindent\textbf{12.} Let $f'(x) = \frac{192x^3}{2 + \sin^4(\pi x)}$ for all $x \in \mathbb{R}$ with $f\left(\frac{1}{2}\right) = 0$.\\
     If $m \leq \int_{1/2}^{1} f(x) dx \leq M$, then the possible values of $m$ and $M$ are $\hfill \textbf{(JEE Adv. 2015)}$
     \begin{multicols}{2}
     \begin{enumerate}
             \item[(a)] $m = 13, M = 24$
	         \item[(b)] $m = \frac{1}{4}, M = \frac{1}{2}$
		     \columnbreak
		         \item[(c)] $m = -11, M = 0$
			     \item[(d)] $m = 1, M = 12$
			     \end{enumerate}
			     \end{multicols}  
			     \textbf{13.} Let
			     \[
			     f(x) = \lim_{n \to \infty} \left( \frac{n^n (x+n) \left(x + \frac{n}{2}\right) \cdots \left(x + \frac{n}{n}\right)}{n! \left(x^2 + n^2\right) \left(x^2 + \frac{n^2}{4}\right) \cdots \left(x^2 + \frac{n^2}{n^2}\right)} \right)^{\frac{x}{n}},
			     \]
			     Then  \text{ for all } $x > 0$. $\hfill \textbf{(JEE Adv. 2015)}$
			     \begin{multicols}{2}
			     \begin{enumerate}
			         \item[(a)] $f\left(\frac{1}{2}\right) \geq f(1)$
				     \item[(b)] $f\left(\frac{1}{3}\right) \leq f\left(\frac{2}{3}\right)$
				         \item[(c)] $f'(2) \leq 0$
					     \item[(d)] $\frac{f'(3)}{f(3)} \geq \frac{f'(2)}{f(2)}$
					     \end{enumerate}
					     \end{multicols}
					     \textbf{14} Let $f: \mathbb{R} \to (0,1)$ be a continuous function. Then, which of the following function(s) has (have) the value zero at some point in the interval (0,1)?  $\hfill \textbf{(JEE Adv 2016)}$
					     \begin{multicols}{2}
					     \begin{enumerate}
					         \item[(a)] $x^9 - f(x)$
						     \item[(b)] $x - \int_{0}^{\frac{\pi}{2} - x} f(t) \cos t \, dt$
						         \item[(c)] $e^x - \int_{0}^{x} f(t) \sin t \, dt$
							     \item[(d)] $f(x) + \int_{0}^{\frac{\pi}{2}} f(t) \sin t \, dt$
							     \end{enumerate}
							     \end{multicols}


							     \textbf{15.} If $g(x) = \int_{\sin x}^{\sin(2x)} \sin^{-1}(t) \, dt$, then
							     $\hfill \textbf{(JEE Adv 2017)}$
							     \begin{multicols}{2}
							     \begin{itemize}
							         \item[(a)] $g'\left(\frac{\pi}{2}\right) = -2\pi$
								     \item[(b)] $g'\left(-\frac{\pi}{2}\right) = 2\pi$
								         \item[(c)] $g'\left(\frac{\pi}{2}\right) = 2\pi$
									     \item[(d)] $g'\left(-\frac{\pi}{2}\right) = -2\pi$
									     \end{itemize}
									     \end{multicols}
									     \textbf{16.} If the line $sx = \alpha$ divides the area of region
									     \[
									     R = \left\{(x, y) \in \mathbb{R}^2 : x^3 \leq y \leq x, 0 \leq x \leq 1\right\}
									     \]
									     into two equal parts, then $\hfill \textbf{(JEE Adv. 2017)}$
									     \begin{multicols}{2}
									     \begin{itemize}
									         \item[(a)] $0 < \alpha \leq \frac{1}{2}$
										     \item[(b)] $\frac{1}{2} < \alpha < 1$
										         \item[(c)] $2\alpha^4 - 4\alpha^2 + 1 = 0$
											     \item[(d)] $\alpha^4 + 4\alpha^2 - 1 = 0$
											     \end{itemize}
											     \end{multicols}
											     \textbf{17.} If $I = \sum_{k=1}^{98} \int_{k}^{k+1} \frac{k+1}{x(k+1)} \, dx$, then $\hfill \textbf{(JEE Adv 2017)}$
											     \begin{multicols}{2}
											     \begin{itemize}
											         \item[(a)] $1 > \log_e 99$  \item[(b)] $1 < \log_e 99$
												     \item[(c)] $1 < \frac{49}{50}$  \item[(d)] $1 > \frac{49}{50}$
												     \end{itemize}
												     \end{multicols}
												     \textbf{18.} For \(a \in \mathbb{R}, |a| \geq 1\), let $\hfill \textbf{(JEE Adv 2019)}$

												     \[
												     \lim_{{n \to \infty}}
												         \frac{
													         1 + \sqrt{2} + \cdots + \sqrt{n}
														     }{
														             (an+1)^2 + (an+2)^2 + \cdots + (an+n)^2
															         } = 54
																 \]
																 Then the possible value(s) of  are:
																 \begin{multicols}{2}
																 \begin{itemize}
																     \item[(a)] $9$
																         \item[(b)] $7$
																	     \item[(c)] $6$
																	         \item[(d)] $8$
																		 \end{itemize}
																		 \end{multicols}
																		 \section*{E Subjective Problems}

																		 1. Find the area bounded by the curve \( x^2 = 4y \) and the straight line \( x = 4y - 2 \). $\hfill \textbf{(1981 - 4 Marks)}$

																		 2. Show that:
																		 \[
																		 \lim_{n \to \infty} \left(1 + \frac{1}{n}\right)^n = e \quad \text{and} \quad \lim_{n \to \infty} \left(1 + \frac{1}{n}\right)^{n^2} = e^{1981}
																		 \]
																		 $\hfill \textbf{(1981 - 4 Marks)}$

																		 3. Show that:
																		 \[
																		 n \int_0^\pi f(x) \sin x \, dx = \frac{n\pi}{2} \int_0^\pi f(\sin x) \, dx
																		 \]
																		 $\hfill \textbf{(1982 - 2 Marks)}$

																		 4. Find the value of:
																		 \[
																		 p \int_0^p |\sin x| \, dx + q \int_p^{2p} \sin x \, dx \quad \text{where} \quad p > q > 0
																		 \]
																		 $\hfill \textbf{(1982 - 3 Marks)}$
																		 5. For any real \( t \), show that the point with coordinates given by:
																		 \[
																		 x = e^{-t}, \quad y = e^{t} - e^{-t}
																		 \]
																		 lies on the hyperbola defined by:
																		 \[
																		 x^2 - \frac{y^2}{e^2} = 1
																		 \]
																		 Show that there are two points corresponding to \( t=1 \) and \( t=-1 \); find them. Show that the area bounded by this hyperbola and these two points corresponding to \( t=1 \) and \( t=-1 \) is
																		 \[
																		 \left(\frac{\pi}{4}\right)(e^2 + 1)
																		 \]
																		 $\hfill \textbf{ (1982 - 3 Marks)}$



																		 \end{document}
